\section{Accessible quantum natural gradient}

Given a supervised loss $L(\bm\theta)$ with Euclidean gradient $g=\nabla_{\bm\theta}L$, AQNG computes a damped preconditioned direction
\begin{equation}
 d_{\rm A}=(\widehat G_{\rm acc}+\lambda I)^{-1}g,
 \qquad
 \bm\theta\leftarrow\bm\theta-\eta d_{\rm A}.
 \label{eq:aqng_update}
\end{equation}
The full-QNG control uses the full quantum metric instead, following quantum natural-gradient optimization \cite{Stokes2020}. From Eq.~\eqref{eq:projection_identity}, $G_{\rm acc}=(P_BS)^{\mathsf T}(P_BS)$, so the accessible metric admits the explicit Gram factorization $A^{\mathsf T}A$ with $A=P_BS$. This factorization permits primal, dual, or SVD-based linear solves depending on parameter and retained-score dimensions. Metric evaluations are cached for a fixed number of optimization steps as an implementation choice.

\begin{figure*}[t]
\centering
\fbox{%
\begin{minipage}{0.94\textwidth}
\textbf{Pseudo-algorithm: Accessible quantum natural gradient (AQNG)}\\[2pt]
\textbf{Input:} initial parameters $\bm\theta_0$, learning rate $\eta$, damping $\lambda$, metric refresh period $K$, supervised minibatches $\{\mathcal B_t\}$, metric minibatches $\{\mathcal M_t\}$, commuting readout family $F$, optional calibration tangents $\mathcal C$.\\[4pt]
\textbf{1. Optional calibration.} If an aligned readout is used, estimate a rank-matched retained subspace from the calibration tangents $\mathcal C$ and freeze it before supervised training.\\[3pt]
\textbf{2. For $t=0,1,\ldots,T-1$:}
\begin{enumerate}
  \item Estimate the supervised gradient $g_t=\nabla_{\bm\theta}L_{\mathcal B_t}(\bm\theta_t)$.
  \item If $t=0$ or $t\bmod K=0$, estimate outcome probabilities and score vectors on $\mathcal M_t$; build the centered readout matrix $B_t$; compute $J_t=B_t^{\mathsf T}S_t$ and $\Sigma_t=B_t^{\mathsf T}B_t$; then form the accessible metric $\widehat G_{{\rm acc},t}=J_t^{\mathsf T}\Sigma_t^{+}J_t$.
  \item Normalize $\widehat G_{{\rm acc},t}$ and apply the damping and trust-region controls.
  \item Solve $d_t=(\widehat G_{{\rm acc},t}+\lambda I)^{-1}g_t$.
  \item Clip $d_t$ if the Euclidean step norm or the metric trust radius exceeds its cap.
  \item Update $\bm\theta_{t+1}=\bm\theta_t-\eta d_t$.
\end{enumerate}
\textbf{3. Output:} final parameters $\bm\theta_T$.
\end{minipage}}
\caption{Pseudo-algorithm for AQNG. The metric can be refreshed less often than the loss gradient and can use an independently sampled metric minibatch.}
\label{fig:aqng_pseudocode}
\end{figure*}

\subsection{Metric normalization and damping}
Metric scale and damping cannot be varied independently: replacing $G$ by $cG$ while holding $\lambda$ fixed changes $(cG+\lambda I)^{-1}$ in a direction-dependent way unless $\lambda$ is rescaled with $c$. The primary configuration therefore normalizes each metric by its trace before applying a common absolute damping coefficient. Additional controls use no normalization or maximum-eigenvalue normalization and alternative damping conventions based on the mean or maximum eigenvalue. We also impose a Euclidean direction cap and an accessible-metric trust radius. These controls bound the accepted update when the estimated metric is nearly singular; consequently, update magnitude alone is not treated as evidence that a vanishing-gradient problem has been solved.

\subsection{Rank-matched readout designs}
The primary readout comparison keeps rank fixed and changes only orientation. The physical readout is the local computational-basis $Z$ span. A random control draws a subspace of exactly the same rank from the centered outcome-function space. The aligned control is fit from an independent calibration set of tangent scores: the leading rank-matched score directions are estimated on calibration tangents and then frozen before optimization and evaluation. Labels are not used in this calibration step. This cross-fit construction avoids evaluating an aligned subspace on the same tangent sample used to fit it.

The rank-matched comparison is important because increasing the number of retained observables changes both dimension and orientation. By holding rank fixed, physical, random, and aligned AQNG test whether orientation relative to the tangent covariance changes optimization at a fixed readout dimension. A separate readout-order ablation deliberately changes rank by moving from weight-one diagonal Pauli functions to all diagonal strings through weight two.

\subsection{Finite-shot estimation}
In finite-shot mode, all retained diagonal features are derived from the same computational-basis bitstrings. Known outcome support is fixed structurally rather than inferred from observed nonzero counts. For the number-conserving control, the support is the corresponding fixed-Hamming-weight sector. A symmetric Dirichlet pseudocount regularizes empirical probabilities on that known support before constructing score and covariance estimates. This prevents unobserved-but-allowed outcomes from producing artificial support collapse at finite shot count.

The finite-shot Full-QNG control uses an analytic full metric while its loss is sampled. It is therefore an optimizer reference, not a hardware-resource-matched baseline. Shot accounting for AQNG reports calibration and training resources separately.
